\subsection{Financial Markets}\label{sec:financial-markets}
A financial market is an organized venue in which participants trade financial assets. A financial asset is a claim to future value.  An exchange is a market that operates under standardized rules, which specify what may be traded and how. 

Buyers need to find sellers that are willing to sell the asset at a low price, while sellers need to find buyers that are willing to pay a high price. The bid is the highest price at which a participant is willing to buy an asset, and the ask is the lowest price at which a participant is willing to sell it. A transaction occurs when the two sides agree on a price. As a result, the observed price of an asset is not a fixed property but rather the outcome of the most recent transaction. A market is considered liquid if many participants can buy or sell continuously, without affecting the asset price \parencite[3--7]{harrisTradingExchangesMarket2003}. 

The return of an asset measures the relative change in its price over a given time period. The risk of an asset is the uncertainty of its return. It is generally measured with the standard deviation of its return, which is commonly referred to as volatility. The risk-free rate is the return of an investment that carries no risk, such as a short-term government bond. The Sharpe ratio is a measure of the risk-adjusted return of an asset
\begin{equation}
\mathrm{SR} = \frac{\mu - r_f}{\sigma},
\label{eq:std-sharpe-ratio}
\end{equation}
where $\mu$ is the mean return over the period, $r_f$ is the risk-free rate and $\sigma$ is the standard deviation of the returns. $\mathrm{SR}$ is commonly annualized.

\textcite[33]{hullOptionsFuturesOther2022} differentiates three main categories of traders. Hedgers trade to reduce the risk they are exposed to from future price movements. Speculators trade to take a position on the future direction of the price. Lastly, arbitrageurs combine long and short positions in two or more instruments, so that a profit is guaranteed regardless of the price movements. A participant who holds a long position benefits from a price increase, and a participant who holds a short position benefits from a price decrease.

Assets are grouped into asset classes according to their characteristics. Equity and commodities are two asset classes that are relevant for this thesis. \textbf{Equity} represents ownership in a company. The holder of equity is entitled to a share of the company's profits. The value of equity therefore depends on the company's future profits. Shares are traded on stock exchanges, financial markets that specialize in equity. \textbf{Commodities} are physical goods that are produced and consumed. The supply of a commodity is determined by production and by the inventories that are held. Demand is generally inelastic, since the good is essential and has few substitutes in the short run. As a result, commodity prices are sensitive to physical conditions such as weather, storage levels and available transport capacity.

As stated above, equity and commodities differ in how their prices are formed. They also differ in what a trade eventually leads to. A trade in equity is settled in money, while a trade in a commodity can end in the delivery of the physical good, which the buyer has to receive and store. Lastly, the two differ in what their holders are exposed to. A holder of equity is affected only by the price of the shares, whereas a gas producer is affected by the price of gas and also by how much of it is sold, which changes with the weather and with economic conditions \parencite{gemanCommoditiesCommodityDerivatives2006}. The European natural gas market is the commodity market studied in this thesis.